\begin{center}\begin{tikzpicture}[font=\small,>=stealth,line join=round]
\draw[gray!40,rounded corners=3pt] (0,0) rectangle (14.2,8.3);\node at (7.1,7.700000000000001) {奇乘偶：一个负号留下，整体仍为奇函数};
\begin{scope}[shift={(4.1,3.9499999999999997)},x=1.875cm,y=0.8846153846153845cm]\draw[->,gray] (-1.6,0)--(1.6,0) node[right]{$x$};\draw[->,gray] (0,-2.6)--(0,2.6) node[above]{$y$};\draw[blue,thick,solid,domain=-1.3:1.3,samples=180] plot(\x,{(\x)^3});\fill[orange!85!black] (-1,-1) circle(2.2pt);\fill[orange!85!black] (1,1) circle(2.2pt);\end{scope}\node[anchor=west,align=left,text width=5.0cm] at (8.4,4.0) {$x\cdot x^2=x^3$\\[10pt]$x$ 左右反号\\[10pt]$x^2$ 左右同号\\[10pt]整体高度互为相反数};
\node[text width=12.6cm,align=center] at (7.1,.65) {乘除规律都在共同定义域上使用；除法还必须保留分母非零条件。};
\end{tikzpicture}\end{center}
